\documentclass[12pt,reqno]{amsart}
\newcommand{\worksheetnumber}{4.3}
\input{../../101preamble.tex}

\DeclareMathOperator{\Res}{Res}
\DeclareMathOperator{\Frac}{Frac}

\title{Worksheet 4.3: Tensor Products, Scalar Extension, and Localization}
\begin{document}
\maketitle
\thispagestyle{firstpage}

Let $R$ be a ring. A \emph{right $R$-module} is an abelian group $(M,+,0_M)$ together with a scalar multiplication $M\times R\to M$, written $(m,r)\mapsto mr$, satisfying the following axioms:
\begin{enumerate}[label=(\roman*)]
\item (Identity) $m1_R=m$ for all $m\in M$,
\item (Associativity) $(mr)s=m(rs)$ for all $r,s\in R$ and all $m\in M$,
\item (Left Distributivity) $(m+n)r=mr+nr$ for all $r\in R$ and all $m,n\in M$, and
\item (Right Distributivity) $m(r+s)=mr+ms$ for all $r,s\in R$ and all $m\in M$.
\end{enumerate}
A homomorphism of right $R$-modules is an additive map $\phi$ with $\phi(mr)=\phi(m)r$. We write $\Mod\text{-}R$ for the category of right $R$-modules, and keep $R\text{-}\Mod$ for left modules. For example, $R$ is a module over itself on either side. Recall from Worksheet 1.2 that when $R$ is commutative, a left module can also be regarded as a right module by setting $mr\coloneqq rm$; without commutativity this formula need not define a right action.

Now let $M$ be a right $R$-module, $N$ a left $R$-module, and $G$ an abelian group. A map $\beta:M\times N\to G$ is \emph{$R$-balanced} if it is additive in each variable and
\[
\beta(mr,n)=\beta(m,rn)\qquad\text{for all }m\in M,\ n\in N,\ r\in R.
\]
The equation lets a scalar pass from the right of the first input to the left of the second input. Notice that we have not written $r\beta(m,n)$: the target is only an abelian group, and it need not have an action of $R$. Matrix multiplication gives a useful model. A row vector is a right module over a matrix ring, a column vector is a left module, and the identity $(mA)n=m(An)$ says exactly that their product is balanced over that matrix ring.

The goal of this worksheet is to extend the tensor product from vector spaces to modules, and to use it to change the ring of scalars. In Worksheet 2.2 we constructed tensor products of vector spaces from bilinear maps, and used splittings to prove that tensoring preserves short exact sequences. Over a ring the construction still works, but as we saw in Worksheet 4.1, the splitting argument does not. We will find that tensor products still respect quotients, while an injection may become the zero map after tensoring; this distinction leads to \emph{flat} modules. We will then use tensor products to extend scalars along a ring homomorphism, and study \emph{localization}, a particularly useful extension of scalars which does preserve exactness.

As in the previous worksheets, all rings have a multiplicative identity but need not be commutative, ring homomorphisms preserve the multiplicative identity, and the letter $\KK$ will denote a field. Since the two inputs to a tensor product over a noncommutative ring play different roles, we will keep careful track of the side on which a ring acts. Towards the end of the worksheet we will assume that $R$ is commutative, and we will say so explicitly.

The \emph{tensor product} $M\otimes_RN$ is an abelian group together with a balanced map $\tau:M\times N\to M\otimes_RN$, written $\tau(m,n)=m\otimes n$, with the following universal property: for every abelian group $G$ and every balanced map $\beta:M\times N\to G$, there is a unique homomorphism of abelian groups $\widetilde\beta:M\otimes_RN\to G$ with $\beta=\widetilde\beta\tau$. In terms of diagrams, we require a unique dashed arrow making
\[
\begin{tikzcd}[row sep=4em,column sep=4em]
M\times N \arrow[r,"\tau"] \arrow[dr,"\beta"'] & M\otimes_RN \arrow[d,dashed,"\widetilde\beta"]\\
& G
\end{tikzcd}
\]
commute. As in Worksheet 2.2, we construct it by starting with freely chosen symbols and imposing the required relations. Let $F\coloneqq\ZZ^{\langle M\times N\rangle}$ be the free abelian group with basis symbols $[m,n]$, and let $H\subset F$ be the subgroup generated by
\[
\begin{gathered}
[m+m',n]-[m,n]-[m',n],\qquad[m,n+n']-[m,n]-[m,n'],\\
[mr,n]-[m,rn],
\end{gathered}
\]
for all $m,m'\in M$, $n,n'\in N$, and $r\in R$. Define $M\otimes_RN\coloneqq F/H$, and let $m\otimes n$ be the class of $[m,n]$. Notice the difference from the construction for vector spaces: the freely chosen coefficients are now integers, so any further scalar action on the tensor product will have to come from additional structure on its inputs. On the other hand, the argument from Worksheet 2.2 shows without change that the universal property determines the tensor product up to a unique isomorphism compatible with the specified balanced maps: two tensor products each receive a unique map from the other, and uniqueness forces both composites to be identities. We check the universal property in the following exercise.

\hrulefill

\begin{problems}
\item\label{t:balanced} \textbf{Balanced Maps and the Tensor Product.} Throughout parts (b) through (e) let $M$ be a right $R$-module and $N$ a left $R$-module.
\begin{enumerate}
\item Verify that row vectors $\KK^{1\times d}$ form a right $\Mat_d(\KK)$-module and column vectors $\KK^{d\times1}$ form a left $\Mat_d(\KK)$-module. Check that multiplying a row by a column is a balanced map to the additive group of $\KK$.
\item Show that the quotient construction makes $(m,n)\mapsto m\otimes n$ balanced. Deduce that $m\otimes0=0\otimes n=0$ and that $(-m)\otimes n=-(m\otimes n)=m\otimes(-n)$. Explain why every element of $M\otimes_RN$ is a finite sum of pure tensors.
\item Extend a balanced map $\beta:M\times N\to G$ to a homomorphism $F\to G$. Show that it kills each generator of $H$, and hence descends to $\widetilde\beta$. Prove uniqueness using the fact that pure tensors generate $M\otimes_RN$. Conversely, show that composing any homomorphism out of $M\otimes_RN$ with $\tau$ gives a balanced map.
\item For a homomorphism of right modules $a:M\to M'$ and a homomorphism of left modules $b:N\to N'$, construct $a\otimes b:M\otimes_RN\to M'\otimes_RN'$ with $(a\otimes b)(m\otimes n)=a(m)\otimes b(n)$. Verify the identity and composition laws, and check that $a\otimes b$ is additive in $a$ and in $b$.
\item Construct isomorphisms $R\otimes_RN\cong N$ and $M\otimes_RR\cong M$ using multiplication. Write down inverse maps, and use the balancing relation to check both composites.
\end{enumerate}
\end{problems}

\hrulefill

What extra structure can a tensor product carry? First suppose $R$ is commutative, so that every $R$-module can be regarded as a module on either side. Then
\[
r(m\otimes n)\coloneqq(rm)\otimes n=m\otimes(rn)
\]
defines an $R$-module structure on $M\otimes_RN$. To see that it is well defined, fix $r\in R$ and consider the map $(m,n)\mapsto(rm)\otimes n$. It is additive in each variable, and it is balanced because, for $s\in R$,
\[
r(ms)=r(sm)=(rs)m=(sr)m=s(rm)=(rm)s,
\]
where the middle equality is exactly where we use commutativity. Hence it induces a homomorphism of $M\otimes_RN$, which is multiplication by $r$, and the module axioms can be checked on pure tensors. With this structure, tensor products represent $R$-bilinear maps to $R$-modules, recovering our construction from Worksheet 2.2 when $R$ is a field.

There is another way to keep scalars, which does not require commutativity. Let $S$ be a second ring. An \emph{$(S,R)$-bimodule} is an abelian group $B$ with a left $S$-action and a right $R$-action satisfying $(sb)r=s(br)$ for all $s\in S$, $b\in B$, and $r\in R$. Notice that the two actions commute with each other, even when multiplication inside either ring is noncommutative. For example, $R$ itself is an $(R,R)$-bimodule. For a left $R$-module $N$, the formula
\[
s(b\otimes n)\coloneqq(sb)\otimes n
\]
makes $B\otimes_RN$ a left $S$-module. The right $R$-action on $B$ is used up in the balancing relation, while the left $S$-action remains available on the result. This is the structure which will allow us to extend scalars along a ring homomorphism. We check both constructions in the following exercise.

\hrulefill

\begin{problems}
\item\label{t:actions} \textbf{Scalar Actions on Tensor Products.}
\begin{enumerate}
\item Suppose $R$ is commutative. For an $R$-module $P$, prove that the homomorphism induced by an $R$-bilinear map $M\times N\to P$ is $R$-linear. Construct a symmetry isomorphism $M\otimes_RN\cong N\otimes_RM$ by interchanging the factors.
\item Now let $B$ be an $(S,R)$-bimodule and $N$ a left $R$-module. For fixed $s\in S$, prove that $(b,n)\mapsto(sb)\otimes n$ is balanced. Identify where you use $(sb)r=s(br)$, and verify that the induced maps give the asserted left $S$-action.
\item If $a:B\to B'$ respects both actions and $c:N\to N'$ is $R$-linear, prove that $a\otimes c$ is $S$-linear. In particular, explain why $B\otimes_R-$ is a functor $R\text{-}\Mod\to S\text{-}\Mod$.
\end{enumerate}
\end{problems}

\hrulefill

We next ask how tensoring interacts with quotients. Fix a left $R$-module $N$ and a homomorphism of right modules $a:M'\to M$. Write $C\coloneqq\coker(a)=M/\img(a)$, and let $q:M\to C$ be the quotient map. We will see that there is an isomorphism of abelian groups
\[
C\otimes_RN\cong\frac{M\otimes_RN}{\img(a\otimes\Id_N)}.
\]
To understand why, consider homomorphisms from the right-hand side to an abelian group $G$. Such homomorphisms correspond to balanced maps on $M\times N$ which vanish on every pair $(a(m'),n)$, and these are exactly the balanced maps which descend in the first input to $C\times N$. Thus both sides have the same maps out. Compare this with the result from Worksheet 2.2 that left adjoints preserve cokernels; here we will turn the explanation into explicit inverse maps.

\begin{theorem}[Right Exactness of Tensor Products]
If $M'\xrightarrow{a}M\xrightarrow{q}M''\to0$ is an exact sequence of right $R$-modules and $N$ is a left $R$-module, then
\[
M'\otimes_RN\xrightarrow{\ a\otimes\Id_N\ }M\otimes_RN\xrightarrow{\ q\otimes\Id_N\ }M''\otimes_RN\longrightarrow0
\]
is an exact sequence of abelian groups. The corresponding statement holds in the second variable, for an exact sequence of left $R$-modules and a fixed right $R$-module.
\end{theorem}

The surjectivity assertion has a simple explanation: lift the first factor of each pure tensor, and add the lifts. Exactness in the middle is subtler, because a tensor which becomes zero need not have an expression in which each individual term becomes zero. The quotient construction handles all of these relations at once. Notice that the theorem makes no claim of injectivity at the left end. For vector spaces, that additional assertion used the existence of a splitting, and the same argument from Worksheet 2.2 shows that tensoring a \emph{split} short exact sequence of modules still gives a split short exact sequence. In general there is no splitting to use. We prove the theorem in the following exercise.

\hrulefill

\begin{problems}
\item\label{t:rightexact} \textbf{Tensor Products and Cokernels.} Throughout parts (a) through (c) use the notation preceding the theorem.
\begin{enumerate}
\item Let $D\coloneqq(M\otimes_RN)/\img(a\otimes\Id_N)$. Show that $([m],n)\mapsto[m\otimes n]$ is well defined on $C\times N$, additive in both variables, and balanced. Construct the induced homomorphism $C\otimes_RN\to D$.
\item Show that $q\otimes\Id_N$ kills $\img(a\otimes\Id_N)$, and hence induces a homomorphism $D\to C\otimes_RN$. Check that the two homomorphisms are inverse to one another on pure tensors.
\item For the exact sequence in the theorem, identify $M''$ with $\coker(a)$ using the first isomorphism theorem. Deduce exactness at $M\otimes_RN$. Prove that $q\otimes\Id_N$ is surjective by lifting finite sums of pure tensors.
\item Repeat the argument in the second variable: for a homomorphism of left modules $b:N'\to N$, prove that $M\otimes_R\coker(b)\cong\coker(\Id_M\otimes b)$. Deduce the second assertion of the theorem.
\end{enumerate}
\end{problems}

\hrulefill

For a first calculation, suppose $R$ is commutative and let $I\subset R$ be an ideal. For an $R$-module $N$, write $IN$ for the submodule of finite sums $\sum_ja_jn_j$ with $a_j\in I$ and $n_j\in N$, as in Worksheet 4.1. We will see that tensoring with $R/I$ imposes the requirement that every element of $I$ act by zero:
\[
(R/I)\otimes_RN\cong N/IN,\qquad\overline r\otimes n\longmapsto rn+IN,
\]
with inverse $n+IN\mapsto\overline1\otimes n$. Thus changing scalars from $R$ to $R/I$ need not preserve the underlying abelian group: it removes the part generated by the action of $I$.

Over $\ZZ$, this becomes $(\ZZ/m\ZZ)\otimes_{\ZZ}N\cong N/mN$. In particular, a tensor product of two finite cyclic groups can be computed by imposing two integer relations on a single generator. For example, $(\ZZ/4\ZZ)\otimes_{\ZZ}(\ZZ/6\ZZ)\cong(\ZZ/6\ZZ)/4(\ZZ/6\ZZ)$, and since $4\cdot\overline1=\overline4$ generates the subgroup $\{\overline0,\overline2,\overline4\}$, this is $\ZZ/2\ZZ$, generated by $\overline1\otimes\overline1$. These calculations also produce our first failure of injectivity. Multiplication by $n$ is injective on $\ZZ$ for $n\geq2$, but it is zero on $\ZZ/n\ZZ$, and after identifying $\ZZ\otimes_{\ZZ}(\ZZ/n\ZZ)$ with $\ZZ/n\ZZ$, this zero map is exactly the result of tensoring the injection with $\ZZ/n\ZZ$. We compute some examples in the following exercise.

\hrulefill

\begin{problems}
\item\label{t:cyclic} \textbf{Quotients, Cyclic Groups, and Failure of Injectivity.} Throughout part (a) let $R$ be commutative, let $I\subset R$ be an ideal, and let $N$ be an $R$-module.
\begin{enumerate}
\item Check that $(\overline r,n)\mapsto rn+IN$ is well defined and balanced. Construct both maps in the displayed isomorphism, and prove that they are $R$-linear and inverse to one another. For the inverse, explain why every element of $IN$ maps to zero.
\item For positive integers $m$ and $n$, identify $(\ZZ/m\ZZ)\otimes_{\ZZ}(\ZZ/n\ZZ)$ with $\ZZ/(m\ZZ+n\ZZ)$. Using the description of the ideals of $\ZZ$ from Worksheet 4.2, show that $m\ZZ+n\ZZ=d\ZZ$, where $d=\gcd(m,n)$, and conclude that the tensor product is $\ZZ/d\ZZ$. Which pure tensor corresponds to $\overline1$?
\item Compute $(\ZZ/6\ZZ)\otimes_{\ZZ}(\ZZ/15\ZZ)$ and $(\ZZ/4\ZZ)\otimes_{\ZZ}(\ZZ/9\ZZ)$, including the image of an arbitrary pure tensor $\overline a\otimes\overline b$ under each identification.
\item Let $n\geq2$, and tensor the short exact sequence $0\to\ZZ\xrightarrow{\times n}\ZZ\to\ZZ/n\ZZ\to0$ with $\ZZ/n\ZZ$. Identify all three resulting groups and both maps. Verify right exactness, and determine the kernel of the first map. Why can no zero be placed at the left end of the tensored sequence while keeping it exact?
\item For $R=\KK[x]$, compute $(R/\langle x\rangle)\otimes_R(R/\langle x^2\rangle)$ and $(R/\langle x\rangle)\otimes_R(R/\langle x-1\rangle)$. Use part (a) to explain your answers without choosing a presentation of either module as an abelian group.
\end{enumerate}
\end{problems}

\hrulefill

A left $R$-module $N$ is \emph{flat} if the functor $-\otimes_RN$ takes every short exact sequence of right $R$-modules to a short exact sequence of abelian groups. Since tensoring is already right exact, flatness asks precisely that tensoring with $N$ preserve injections. Notice that the handedness in this definition matters: the fixed module carries a left action, so the variable modules carry right actions. For a commutative ring, the symmetry isomorphism of Problem~\ref{t:actions}(a) lets us suppress this distinction.

It is natural to ask which modules are flat. Free modules are, because tensoring with a free module gives a direct sum of copies of the original abelian group. More generally, tensoring commutes with arbitrary direct sums, in either variable: for a family $(N_j)_{j\in J}$ of left modules,
\[
M\otimes_R\Bigl(\bigoplus_jN_j\Bigr)\cong\bigoplus_j(M\otimes_RN_j),\qquad m\otimes(n_j)_j\longmapsto(m\otimes n_j)_j.
\]
The finite-support condition is what makes this formula possible: only finitely many coordinates appear in each pure tensor, and hence in each finite sum of pure tensors. This is the same reason that a family of maps out of a direct sum could be combined in Worksheet 4.1. Notice also that a direct sum of injective homomorphisms is injective, since an element of a direct sum is zero exactly when all of its coordinates are zero; the same coordinatewise argument shows that a direct sum of short exact sequences is again short exact.

Recall from Worksheet 4.2 that projective modules are exactly the direct summands of free modules. We will see that a direct summand of a flat module is flat, so every projective module is flat. These properties are nevertheless different: projectivity supplies lifts of maps, while flatness preserves injections after tensoring. Localization will give a concrete flat module which is not projective. First we establish the implications which follow from free modules and splittings in the following exercise.

\hrulefill

\begin{problems}
\item\label{t:flat} \textbf{Direct Sums, Free Modules, and Flatness.} Throughout this exercise let $M$ be a right $R$-module.
\begin{enumerate}
\item Construct the displayed isomorphism and its inverse, using the coordinate inclusions. Verify the balancing relation, and check both composites on pure tensors. Prove the corresponding isomorphism $\bigl(\bigoplus_jM_j\bigr)\otimes_RN\cong\bigoplus_j(M_j\otimes_RN)$ for a family of right modules $(M_j)_{j\in J}$.
\item Let $R^{\langle J\rangle}$ be a free left module. Use part (a) and the isomorphism $M\otimes_RR\cong M$ to identify $M\otimes_RR^{\langle J\rangle}$ with $\bigoplus_{j\in J}M$. Verify that tensoring a homomorphism with $R^{\langle J\rangle}$ gives its coordinatewise direct sum. Deduce that every free left module is flat, including free modules of infinite rank and the zero module.
\item Suppose $F=N\oplus Q$ is flat. For an injective homomorphism of right modules $a:M'\to M$, identify $a\otimes\Id_F$ with the direct sum of $a\otimes\Id_N$ and $a\otimes\Id_Q$. Deduce that both summands are injective, and hence that $N$ and $Q$ are flat.
\item Combine the direct summand characterization of projective modules from Worksheet 4.2 with parts (b) and (c) to prove that every projective left module is flat. Which example in Problem~\ref{t:cyclic} shows that not every module is flat?
\end{enumerate}
\end{problems}

\hrulefill

Let $f:R\to S$ be a ring homomorphism. A left $S$-module $P$ becomes a left $R$-module by $rp\coloneqq f(r)p$. This is called \emph{restriction of scalars}, and we write $\Res_f(P)$ for the resulting $R$-module. It changes which scalars we allow, but keeps the underlying abelian group. In the other direction, $S$ is an $(S,R)$-bimodule, with its usual left multiplication and the right action $sr\coloneqq sf(r)$; the bimodule axioms follow from associativity in $S$ and the fact that $f$ is a ring homomorphism preserving the multiplicative identity. Thus a left $R$-module $N$ gives a left $S$-module
\[
S\otimes_RN,
\]
called its \emph{extension of scalars}. The map $\eta_N:N\to\Res_f(S\otimes_RN)$, $n\mapsto1\otimes n$, is $R$-linear. Notice that it need not be injective, as the example $(R/I)\otimes_RN\cong N/IN$ already shows.

The universal property of extension of scalars says that an $R$-linear map $g:N\to\Res_f(P)$ extends uniquely to an $S$-linear map $S\otimes_RN\to P$, given by $s\otimes n\mapsto sg(n)$, and that restricting this map along $\eta_N$ recovers $g$. We therefore expect a natural bijection
\[
\Hom_S(S\otimes_RN,P)\cong\Hom_R\bigl(N,\Res_f(P)\bigr),
\]
that is, an adjunction in the sense of Worksheet 2.2. Unlike restriction, extension of scalars generally changes the abelian group. For $\ZZ\to\QQ$ it makes division by nonzero integers possible, and for $R\to R/I$ it makes $I$ act by zero. The adjunction describes both operations by their maps into modules over the new ring. We check this in the following exercise.

\hrulefill

\begin{problems}
\item\label{t:scalars} \textbf{Extension and Restriction of Scalars.} Throughout this exercise let $f:R\to S$ be a ring homomorphism.
\begin{enumerate}
\item Check that restriction of scalars is a functor which does not change the underlying functions of homomorphisms, and use Problem~\ref{t:actions} to explain why extension of scalars is a functor. Prove that $\eta_N$ is $R$-linear. When $R$ is commutative and $S=R/I$, identify $\eta_N$ with the quotient map $N\to N/IN$ using Problem~\ref{t:cyclic}(a).
\item For an $R$-linear map $g:N\to\Res_f(P)$, check that $(s,n)\mapsto sg(n)$ is balanced. Construct the induced homomorphism $\widetilde g:S\otimes_RN\to P$, and prove that it is $S$-linear. Conversely, for an $S$-linear map $F:S\otimes_RN\to P$, prove that $n\mapsto F(1\otimes n)$ is $R$-linear.
\item Show that these two assignments are inverse to one another. Write $\Theta_{N,P}(F)(n)\coloneqq F(1\otimes n)$. If $a:N'\to N$ is $R$-linear and $b:P\to P'$ is $S$-linear, verify the naturality equation
\[
\Theta_{N',P'}\bigl(bF(\Id_S\otimes a)\bigr)=\Res_f(b)\,\Theta_{N,P}(F)\,a.
\]
Conclude that extension of scalars is left adjoint to restriction of scalars.
\item Show that restriction of scalars is exact and that extension of scalars is right exact. Use $\ZZ\to\ZZ/n\ZZ$ to show that extension of scalars need not be exact. How does right exactness relate to the result from Worksheet 2.2 that left adjoints preserve cokernels?
\end{enumerate}
\end{problems}

\hrulefill

Extension of scalars is a special case of a tensor--Hom adjunction involving two rings. Let $B$ be an $(S,R)$-bimodule and let $P$ be a left $S$-module. The abelian group $\Hom_S(B,P)$ becomes a left $R$-module by
\[
(r\cdot h)(b)\coloneqq h(br).
\]
Notice that the scalar acts by precomposing with the right action on $B$, so we do not need an action of $R$ on $P$. Compare this with the coinduced modules of Worksheet 4.2. For a left $R$-module $N$, the adjunction is
\[
\Hom_S(B\otimes_RN,P)\cong\Hom_R\bigl(N,\Hom_S(B,P)\bigr).
\]
The maps implementing it are the familiar ones from Worksheet 2.2: an $S$-linear map $F$ corresponds to $\widehat F$ with $\widehat F(n)(b)\coloneqq F(b\otimes n)$, and an $R$-linear map $g$ corresponds to the map with $b\otimes n\mapsto g(n)(b)$. The order of the letters differs from the formula for vector spaces, because the right $R$-action belongs to $B$. It is worth seeing exactly how the balancing relation becomes $R$-linearity. Using $(br)\otimes n=b\otimes(rn)$, we find
\[
\widehat F(rn)(b)=F\bigl(b\otimes rn\bigr)=F\bigl(br\otimes n\bigr)=\widehat F(n)(br)=\bigl(r\cdot\widehat F(n)\bigr)(b),
\]
which says precisely that $\widehat F(rn)=r\cdot\widehat F(n)$. We check the remaining details in the following exercise.

\hrulefill

\begin{problems}
\item\label{t:adjunction} \textbf{Tensor--Hom with Compatible Actions.} Throughout this exercise let $B$ be an $(S,R)$-bimodule, let $N$ be a left $R$-module, and let $P$ be a left $S$-module.
\begin{enumerate}
\item For $h\in\Hom_S(B,P)$ and $r\in R$, prove that $b\mapsto h(br)$ is $S$-linear. Check the left $R$-module axioms for $\Hom_S(B,P)$. In particular, compute $r_1\cdot(r_2\cdot h)$ and compare it with $(r_1r_2)\cdot h$.
\item Given an $S$-linear map $F:B\otimes_RN\to P$, prove that each $\widehat F(n)$ is $S$-linear in $b$, and that $n\mapsto\widehat F(n)$ is additive. Together with the computation above, deduce that it is $R$-linear.
\item Given an $R$-linear map $g:N\to\Hom_S(B,P)$, prove that $(b,n)\mapsto g(n)(b)$ is additive in each variable and balanced, and that the induced homomorphism out of the tensor product is $S$-linear. Show that the two constructions are inverse homomorphisms of abelian groups.
\item For an $R$-linear map $a:N'\to N$ and an $S$-linear map $c:P\to P'$, verify that
\[
\widehat{cF(\Id_B\otimes a)}(n')(b)=c\bigl(\widehat F(a(n'))(b)\bigr).
\]
Show also that postcomposition with $c$ is an $R$-linear map between the Hom modules. Deduce naturality in both variables, and state which functor is left adjoint to which.
\item Take $B=S$, with the bimodule structure described before Problem~\ref{t:scalars}. Prove that evaluation at $1$ gives an $R$-linear isomorphism $\Hom_S(S,P)\cong\Res_f(P)$, and recover the adjunction between extension and restriction of scalars.
\end{enumerate}
\end{problems}

\hrulefill

A concrete extension of scalars already contains a useful surprise. Regard $\CC$ as a two-dimensional vector space over $\RR$, and form $\CC\otimes_{\RR}\CC$. By Worksheet 3.1 this is a commutative algebra, with multiplication $(z\otimes w)(z'\otimes w')=zz'\otimes ww'$. We make it a $\CC$-algebra using the first factor, so that $c(z\otimes w)\coloneqq(cz)\otimes w$. Writing the second factor in the real basis $(1,i)$ shows that every tensor is a complex linear combination of $1\otimes1$ and $1\otimes i$.

There are two natural ways to multiply the two factors together: $z\otimes w\mapsto zw$ and $z\otimes w\mapsto z\overline w$. Both are balanced over $\RR$, and together they give
\[
\Phi:\CC\otimes_{\RR}\CC\longrightarrow\CC\times\CC,\qquad z\otimes w\longmapsto(zw,z\overline w),
\]
where the target has coordinatewise multiplication and the diagonal action $c(a,b)\coloneqq(ca,cb)$. We will see that $\Phi$ is an isomorphism of $\CC$-algebras. For instance, $\Phi(1\otimes1)=(1,1)$ and $\Phi(1\otimes i)=(i,-i)$, and these are linearly independent over $\CC$. Since $1\otimes1$ and $1\otimes i$ span $\CC\otimes_{\RR}\CC$, once we know that $\Phi$ is $\CC$-linear it follows that $\Phi$ is bijective, and that $\CC\otimes_{\RR}\CC$ has dimension two over $\CC$, and four over $\RR$. In particular, extending scalars can produce a product of two fields, even though both of the original algebras were fields. The two coordinates record the two ways in which the second factor can enter the calculation: directly, and through complex conjugation, which is the nontrivial element of $\Aut_{\RR}(\CC)$ from Worksheet 1.1. We check this in the following exercise.

\hrulefill

\begin{problems}
\item\label{t:complex} \textbf{Complex Scalars Extended from the Reals.}
\begin{enumerate}
\item Justify the multiplication on $\CC\otimes_{\RR}\CC$ using the tensor product of algebras from Worksheet 3.1. Check that both coordinates of $\Phi$ come from maps which are balanced over $\RR$, and prove that $\Phi$ is $\CC$-linear, multiplicative, and preserves the multiplicative identity.
\item Solve $(x+iy,x-iy)=(a,b)$ for $x,y\in\CC$. Use your answer to prove that
\[
(a,b)\longmapsto\frac{a+b}{2}\otimes1+\frac{i(b-a)}{2}\otimes i
\]
is inverse to $\Phi$, by checking both composites on the indicated generators.
\item Find the preimages of $(1,0)$ and $(0,1)$ under $\Phi$. Show that they are nonzero idempotents whose product is zero and whose sum is $1\otimes1$. Conclude that $\CC\otimes_{\RR}\CC$ is not a field.
\item Compare $i\otimes1$ and $1\otimes i$ using $\Phi$. Why does balancing over $\RR$ not identify them? Compare this with $\CC\otimes_{\CC}\CC$.
\end{enumerate}
\end{problems}

\hrulefill

For the rest of this worksheet, $R$ is commutative. We now construct an extension of scalars which preserves injections as well as quotients. The familiar passage from $\ZZ$ to $\QQ$ makes every nonzero integer invertible. More generally, we can specify which elements we want to invert. A subset $U\subset R$ is \emph{multiplicatively closed} if $1\in U$, and $uv\in U$ whenever $u,v\in U$. We allow $0\in U$, although we will see that this forces the resulting ring to be the zero ring.

The \emph{localization} $U^{-1}R$ consists of fractions $r/u$ with $r\in R$ and $u\in U$, subject to the relation
\[
\frac ru=\frac{r'}{u'}\qquad\Longleftrightarrow\qquad t(u'r-ur')=0\text{ for some }t\in U.
\]
Addition and multiplication are given by the familiar formulas
\[
\frac ru+\frac sv\coloneqq\frac{vr+us}{uv}\qquad\text{and}\qquad\frac ru\cdot\frac sv\coloneqq\frac{rs}{uv}.
\]
The relation is reflexive and symmetric, and the multiplier is what makes it transitive. Suppose $t(u'r-ur')=0$ and $t'(u''r'-u'r'')=0$. Multiplying the first equation by $t'u''$ and the second by $tu$ and adding, the terms involving $r'$ cancel, leaving $tt'u'(u''r-ur'')=0$, so the multiplier $tt'u'\in U$ shows that $r/u=r''/u''$. The extra multiplier $t$ also accounts for zero divisors, in the sense of Worksheet 1.1. If an element of $U$ kills $r$, then making that element invertible must also make $r$ zero, and ordinary cross multiplication would not impose all of the required identifications. An \emph{integral domain} is a nonzero commutative ring with no zero divisors, that is, in which $ab=0$ forces $a=0$ or $b=0$. When $R$ is an integral domain and $0\notin U$, the extra multiplier can be cancelled, and the relation becomes the familiar equality of fractions.

The map $\eta:R\to U^{-1}R$, $r\mapsto r/1$, sends each $u\in U$ to a unit, with inverse $1/u$. We will see that it has the following universal property: for every ring homomorphism $f:R\to A$ to a commutative ring with $f(U)\subset A^\times$, there is a unique ring homomorphism $\widetilde f:U^{-1}R\to A$ with $\widetilde f\eta=f$. Notice that its values are already prescribed:
\[
\widetilde f(r/u)=f(r)f(u)^{-1}.
\]
The equivalence relation is what makes this prescription well defined, including when $R$ has zero divisors. We construct the localization in the following exercise.

\hrulefill

\begin{problems}
\item\label{t:localring} \textbf{Constructing a Localization.} Throughout parts (a) through (c) let $U\subset R$ be multiplicatively closed.
\begin{enumerate}
\item Show that changing a representative does not change the result of addition or multiplication. Verify that these operations make $U^{-1}R$ a commutative ring, with zero $0/1$, multiplicative identity $1/1$, and additive inverse $(-r)/u$. Hint: common denominators reduce each ring axiom to an identity in $R$.
\item Prove that $\eta$ is a ring homomorphism and that $u/1$ has inverse $1/u$ for $u\in U$. Establish the universal property by checking that the formula for $\widetilde f$ is independent of representatives and preserves sums and products. Prove uniqueness by writing $r/u=(r/1)(u/1)^{-1}$.
\item Show that $r/1=0$ exactly when $tr=0$ for some $t\in U$. Deduce that $\eta$ is injective exactly when no element of $U$ is a zero divisor. Show that $U^{-1}R=0$ if and only if $0\in U$, and apply this when $U$ contains a nilpotent element.
\item Let $R=\ZZ$. Identify $U^{-1}\ZZ$ first for $U=\ZZ\smallsetminus\{0\}$, and then for $U=\{1,2,2^2,\ldots\}$; we write the second ring as $\ZZ[1/2]$. For $R=\ZZ/6\ZZ$ and $U$ the powers of $\overline2$, prove that the localization is isomorphic to $\ZZ/3\ZZ$, and find $\ker(\eta)$.
\item Let $R=\KK[x]/\langle x^2\rangle$ be the dual numbers from Worksheet 1.1, and invert the powers of the class of $x$. Identify the result, and compare it with inverting the powers of $x$ in $\KK[x]$. Which feature of the original ring accounts for the difference?
\end{enumerate}
\end{problems}

\hrulefill

The same construction allows fractions with a module element in the numerator. For an $R$-module $M$, define $U^{-1}M$ to consist of fractions $m/u$ with $m\in M$ and $u\in U$, subject to the relation
\[
\frac mu=\frac{m'}{u'}\qquad\Longleftrightarrow\qquad t(u'm-um')=0\text{ for some }t\in U.
\]
Addition uses common denominators, and $U^{-1}R$ acts by $(r/s)(m/u)\coloneqq(rm)/(su)$, so that $U^{-1}M$ is a module over the localized ring. Notice that the map $M\to U^{-1}M$, $m\mapsto m/1$, can lose elements, just as the map on rings can. The precise criterion is
\[
m/u=0\qquad\Longleftrightarrow\qquad tm=0\text{ for some }t\in U.
\]
We will use this criterion repeatedly, rather than cancelling a denominator inside $M$, where it need not act injectively.

Suppose $P$ is a $U^{-1}R$-module and $g:M\to P$ is $R$-linear after restriction of scalars along $\eta$. The only possible extension of $g$ to fractions is $m/u\mapsto(u/1)^{-1}g(m)$. This resembles the adjunction between extension and restriction of scalars, and in fact localization of modules \emph{is} extension of scalars along $\eta$:
\[
U^{-1}R\otimes_RM\cong U^{-1}M,\qquad\frac ru\otimes m\longmapsto\frac{rm}u.
\]
Both descriptions are useful: tensors explain the universal property, while fractions let us check exactness one element at a time. We construct this isomorphism explicitly in the following exercise.

\hrulefill

\begin{problems}
\item\label{t:modulefractions} \textbf{Fractions in Modules and Extension of Scalars.} Throughout this exercise let $U\subset R$ be multiplicatively closed and let $M$ be an $R$-module.
\begin{enumerate}
\item Adapt the proof of transitivity given before Problem~\ref{t:localring} to module fractions. Check that addition and the displayed scalar multiplication are well defined and make $U^{-1}M$ a $U^{-1}R$-module. Prove the criterion for a fraction to be zero, and show that $m/u=(tm)/(tu)$ for every $t\in U$.
\item For an $R$-linear map $a:M\to N$, show that $m/u\mapsto a(m)/u$ is a well-defined $U^{-1}R$-linear map $U^{-1}M\to U^{-1}N$. Check the identity and composition laws, so that localization is a functor.
\item Let $P$ be a $U^{-1}R$-module. Construct the unique $U^{-1}R$-linear extension $U^{-1}M\to P$ of an $R$-linear map $g:M\to P$. Use the multiplier in the equivalence relation to prove that its formula does not depend on representatives.
\item Check that $(r/u,m)\mapsto(rm)/u$ is $R$-balanced, and construct the displayed homomorphism out of $U^{-1}R\otimes_RM$. Prove that it is $U^{-1}R$-linear.
\item Define an inverse by $m/u\mapsto(1/u)\otimes m$. To show that it is well defined, suppose $t(u'm-um')=0$, and express the difference of the two proposed tensors as
\[
\frac1{tuu'}\otimes t(u'm-um').
\]
Check additivity and compatibility with scalars, and then check both composites on fractions and on pure tensors. Verify that the isomorphism commutes with the maps induced by an $R$-linear map $a:M\to N$.
\end{enumerate}
\end{problems}

\hrulefill

Localization preserves exactness, even though the maps $M\to U^{-1}M$ need not be injective. These two assertions concern different maps. Exactness says that an injection \emph{between two modules} remains an injection after both modules have been localized. If $a:M'\to M$ is injective and a fraction $a(m')/u$ is zero, then some $t\in U$ kills $a(m')$. Injectivity of $a$ then says that the same $t$ kills $m'$, so the fraction $m'/u$ was already zero in $U^{-1}M'$. Thus localization preserves injections.

Exactness in the middle uses a related argument. A numerator need not map to zero before localization. However, if its fraction maps to zero, then multiplying its numerator and denominator by a suitable $t\in U$ gives an equal fraction whose numerator does map to zero, and exactness of the original sequence then provides a preimage. This use of an extra multiplier is exactly why the relation defining fractions had to allow for zero divisors.

\begin{theorem}[Exactness of Localization]
Let $R$ be a commutative ring and $U\subset R$ a multiplicatively closed subset. Localization takes every short exact sequence of $R$-modules to a short exact sequence of $U^{-1}R$-modules. Consequently, $U^{-1}R$ is a flat $R$-module.
\end{theorem}

We prove the theorem in the following exercise.

\hrulefill

\begin{problems}
\item\label{t:localexact} \textbf{Exactness and Flatness of Localization.} Throughout parts (a) and (b) let $0\to A\xrightarrow{a}B\xrightarrow{b}C\to0$ be a short exact sequence of $R$-modules.
\begin{enumerate}
\item Prove that $U^{-1}b$ is surjective, by lifting the numerator of a fraction in $U^{-1}C$. Check that the two localized maps compose to zero.
\item Suppose $b(y)/u=0$. Find $t\in U$ with $b(ty)=0$, and then find $x\in A$ with $a(x)=ty$. Show that $x/(tu)$ maps to $y/u$, and deduce exactness at $U^{-1}B$.
\item Use Problem~\ref{t:modulefractions} and the symmetry isomorphism of Problem~\ref{t:actions}(a) to identify the functor $-\otimes_RU^{-1}R$ with localization. Deduce that $U^{-1}R$ is flat, in the sense of Problem~\ref{t:flat}.
\item Localize $0\to\ZZ\xrightarrow{\times6}\ZZ\to\ZZ/6\ZZ\to0$, first at the powers of $2$, and then at all nonzero integers. Identify each localized sequence, including the final quotient. Why can an exact functor such as localization send a nonzero module to zero?
\end{enumerate}
\end{problems}

\hrulefill

A particularly important choice of $U$ is the complement of a prime ideal. An ideal $\fp\subset R$ is \emph{prime} if $\fp\neq R$, and $ab\in\fp$ implies $a\in\fp$ or $b\in\fp$. Equivalently, $R/\fp$ is an integral domain: a product $\overline a\,\overline b$ is zero in $R/\fp$ exactly when $ab\in\fp$, and $\overline a=0$ exactly when $a\in\fp$. Similarly, an ideal $\fp\neq R$ is prime exactly when its complement is closed under multiplication, and this complement contains $1$ since $\fp\neq R$. Every maximal ideal is prime, since by Worksheet 4.1 its quotient is a field. The complement $R\smallsetminus\fp$ is multiplicatively closed, and we write
\[
R_{\fp}\coloneqq(R\smallsetminus\fp)^{-1}R.
\]
The ring $R_{\fp}$ makes every element outside $\fp$ invertible. Its fractions with numerator in $\fp$ form an ideal
\[
\fp R_{\fp}\coloneqq\left\{r/u \;\; \big| \;\; r\in\fp,\ u\notin\fp\right\},
\]
and we will see that every fraction outside this ideal is a unit. Consequently $\fp R_{\fp}$ is the unique maximal ideal of $R_{\fp}$. A commutative ring with exactly one maximal ideal is called a \emph{local ring}.

It is worth being careful with the notation, since it can conceal an important difference. For a prime number $p$, the ring $\ZZ_{(p)}$ is the localization at the prime ideal $\langle p\rangle$, which allows denominators not divisible by $p$; in particular, $p$ is not a unit. The ring $\ZZ[1/p]$ allows powers of $p$ as denominators, and makes $p$ a unit. In general, we write $R[1/f]$ for the localization at the powers of an element $f$, and keep subscripts for the localization at a prime ideal. Both constructions are described by fractions, but they invert different sets.

For an integral domain $D$, localizing at all of its nonzero elements gives a field, called its \emph{field of fractions} and denoted $\Frac(D)$: a nonzero fraction $a/b$ has inverse $b/a$. We will see that the quotient $R_{\fp}/\fp R_{\fp}$ is isomorphic to $\Frac(R/\fp)$, and we call it the \emph{residue field} of the local ring $R_{\fp}$. Notice that we will construct the quotient map explicitly, so this does not require a separate description of the ideals of a localization. We study localization at a prime in the following exercise.

\hrulefill

\begin{problems}
\item\label{t:prime} \textbf{Localization at a Prime and Local Rings.} Throughout parts (a) through (c) let $\fp\subset R$ be a prime ideal.
\begin{enumerate}
\item Show that whether a fraction $r/u\in R_{\fp}$ has numerator in $\fp$ does not depend on the chosen representative. Use the multiplier in the relation defining fractions, together with the fact that all denominators and multipliers lie outside $\fp$.
\item Prove that $\fp R_{\fp}$ is an ideal of $R_{\fp}$ which is not the whole ring. Show that a fraction $r/u$ is a unit if and only if $r\notin\fp$: construct an inverse in one direction, and in the other use the fact that $\fp R_{\fp}$ is not the whole ring. Deduce that every ideal of $R_{\fp}$ other than $R_{\fp}$ itself is contained in $\fp R_{\fp}$, so that $\fp R_{\fp}$ is the unique maximal ideal.
\item Show that $r/u\mapsto\overline r/\overline u$ defines a surjective ring homomorphism $R_{\fp}\to\Frac(R/\fp)$ with kernel $\fp R_{\fp}$. Apply the first isomorphism theorem to identify the residue field.
\item Let $R=\ZZ$ and $\fp=\langle p\rangle$ for a prime number $p$. Describe $\ZZ_{(p)}$, its units, its maximal ideal, and its residue field. Do the same for $\fp=\langle0\rangle$. Compare $\ZZ_{(p)}$ with $\ZZ[1/p]$ by deciding whether $p$, and a different prime $q$, are units in each.
\item Let $R=\KK[x]$ and $\fp=\langle x\rangle$. Describe the allowed denominators, and identify the map to the residue field with evaluation at $0$. Why is $1+x$ a unit in $R_{\fp}$, although it is not a unit in $\KK[x]$?
\end{enumerate}
\end{problems}

\hrulefill

We can now distinguish flatness from projectivity using a familiar abelian group. Since $\QQ$ is a localization of $\ZZ$, the exactness theorem shows that $\QQ$ is flat over $\ZZ$. To show that $\QQ$ is not projective, we use the direct summand characterization from Worksheet 4.2: if $\QQ$ were projective, it would embed as a direct summand of a free abelian group. But as you showed in Worksheet 4.2, every homomorphism $\QQ\to\ZZ$ is zero, since its values must be divisible by every positive integer. Composing with coordinate projections then shows that every homomorphism from $\QQ$ to a free abelian group is zero.

Notice that this argument uses a property of the maps \emph{out of} $\QQ$, while flatness came from the behavior of tensor products. In particular, we do not need to know anything about subgroups of free abelian groups. Together with the projective module $\ZZ/2\ZZ$ over $\ZZ/6\ZZ$ from Worksheet 4.2, which is not free, this separates three properties which coincide for vector spaces: being free, being projective, and being flat. We carry out the argument in the following exercise.

\hrulefill

\begin{problems}
\item\label{t:rationals} \textbf{A Flat Module Which Is Not Projective.}
\begin{enumerate}
\item Identify $\QQ$ with the localization of $\ZZ$ at the nonzero integers, using either the universal property or the construction by fractions. Use Problem~\ref{t:localexact} to prove that $\QQ$ is flat over $\ZZ$.
\item Let $F=\bigoplus_{j\in J}\ZZ$ be a free abelian group, and let $g:\QQ\to F$ be a homomorphism. Compose $g$ with each coordinate projection, and use $\Hom_{\ZZ}(\QQ,\ZZ)=0$ from Worksheet 4.2, to prove that $g=0$. Explain why an element of $F$ whose coordinates all vanish is zero, even when $J$ is infinite.
\item If $\QQ$ were projective, use the direct summand characterization to produce an injective homomorphism $\QQ\to F$ for some free abelian group $F$. Derive a contradiction from part (b), and conclude that the implication ``projective implies flat'' cannot be reversed.
\item Compute $\QQ\otimes_{\ZZ}\ZZ/n\ZZ$ for $n\geq2$, both by localization and using Problem~\ref{t:cyclic}(a). Compare your answer with $\QQ\otimes_{\ZZ}\ZZ\cong\QQ$, and explain why killing a nonzero module does not contradict flatness.
\end{enumerate}
\end{problems}

\hrulefill

\end{document}
